\section{Related work and the gap}
\begin{frame}{Prior work establishes insecurity, poisoning, and human factors}
\begin{itemize}
\item \textbf{Generated-code security}
  \begin{itemize}
  \item \small Pearce et al. (2022); Khoury et al. (2023): insecure implementations and defaults.
  \end{itemize}
\item \textbf{Developer interaction}
  \begin{itemize}
  \item \small Perry et al. (2023); Sandoval et al. (2023): trust, prompting, and verification affect security.
  \end{itemize}
\item \textbf{Adversarial manipulation}
  \begin{itemize}
  \item \small Schuster et al. (2021); Oh et al. (2024): poisoning can steer assistants toward unsafe output.
  \end{itemize}
\end{itemize}
\medskip
\textbf{Gap:} NES implicitly collects an ongoing edit stream and proactively proposes changes. This study connects these mechanisms to security failures.
\source{Sections 1 and 7; references [27, 37, 43--46]}
\note{These are strands of prior work, not claims that all earlier assistants use only local context. Summaries follow the supplied paper's related-work discussion. Selected references are available in the optional backup slides.}
\end{frame}
